\chapter{El personaje potenciado: Alex Etxebarria}

\section{Mecánica de control compartido}

En \emph{Kids on Bikes}, los personajes con habilidades sobrenaturales no son controlados por el Director ni por un único jugador, sino que son \lead{Personajes Potenciados} (\emph{Powered Characters}) cuya voluntad se comparte entre toda la mesa.

\kobfigure{04_aparicion_alex.png}{Alex Etxebarria en la noche: mirada ausente y destellos de energía verde esmeralda.}

Al encontrar a Alex en la Escena~5, el Director reparte cuatro \lead{Aspectos de Alex} entre los jugadores:
\begin{enumerate}
  \item \textbf{El guardián de los recuerdos:} Decide cuándo le viene a Alex un recuerdo fragmentado de su vida antes de la desaparición o de su tiempo cautivo en B\&W.
  \item \textbf{La brújula emocional:} Expresa los miedos, la tristeza o el apego de Alex hacia sus amigos en momentos de tensión.
  \item \textbf{El canalizador del Éter:} Describe las manifestaciones físicas involuntarias del poder (estática en el pelo, zumbido en los oídos, frío repentino).
  \item \textbf{El protector del grupo:} Decide cómo reacciona Alex instintivamente cuando un amigo está a punto de recibir daño.
\end{enumerate}

\section{Los poderes del Éter y su coste}

Para que Alex use sus habilidades de manera intencionada, los jugadores deben invertir sus \lead{Fichas de Adversidad}:

\begin{CajaEter}{Catálogo de habilidades de Alex}
\begin{itemize}
  \item \lead{Telequinesis menor (Coste: 1 ficha):} Levantar o mover objetos de hasta 10~kg, desarmar a un enemigo, atrancar una puerta o accionar una palanca lejana.
  \item \lead{Visión de energía (Coste: 1 ficha):} Percibir el rastro esmeralda del Éter a través de paredes de hormigón, detectar cámaras de seguridad o identificar cables de alta tensión.
  \item \lead{Escudo de resonancia (Coste: 2 fichas):} Desviar una ráfaga de disparos, frenar el impacto de un vehículo o amortiguar una caída mortal de todo el grupo.
  \item \lead{Distorsión temporal (Coste: 3 fichas):} Congelar la percepción del tiempo durante unos segundos, permitiendo a los personajes escapar de una emboscada o resolver una acción antes de que los guardias reaccionen.
\end{itemize}
\end{CajaEter}

\section{Estrés y sobrecarga psíquica}

Alex no puede canalizar energía cósmica ilimitada sin pagar un precio biológico. Cada vez que Alex utiliza un poder, el jugador que lo activa lanza un dado de estrés (1d6):
\begin{itemize}
  \item Con un resultado de \textbf{1 o 2}, Alex absorbe el esfuerzo con fatiga menor (un leve sangrado nasal o mareo).
  \item Con un resultado de \textbf{3 a 5}, se produce una \lead{Fuga de Éter}: las luces del entorno estallan en chispas verdes, las radios emiten estática ensordecedora y el radio de búsqueda de Kessler recibe una alerta de telemetría.
  \item Con un resultado de \textbf{6}, Alex sufre una \lead{Sobrecarga Cósmica}: pierde el conocimiento temporalmente, levita en el aire envuelto en una tormenta electromagnética que arroja a amigos y enemigos al suelo, y queda indefenso durante los siguientes minutos.
\end{itemize}
